\subsection{Emitterschaltung}





\subsection{Kollektorschaltung}
\paragraph{Länge des Koaxialkabels}
Die bestimmte Länge und angegebene Länge haben mit $S=0.14\sigma$ eine sehr geringe Abweichung. Jedoch sei hier angemerkt, dass die Abschätzung der Signalgeschwindigkeit auf 2/3 der Lichtgeschwindigkeit recht ungenau ist, in der Berechnung aber als perfekt angenommen wurde.

\paragraph{Impedanz des Koaxialkabels} Im Versuch wurde die Impedanz des verwendeten Koaxialkabels zu $R=54.60(21)$ bestimmt. Dies entspricht den typischen Werte für Koaxialkabel, die laut Wikipedia \autocite{Koaxialkabel} bei \(50-75 \, \Omega\) liegen.


\subsection{PID-Regler}

\subsubsection{Interner Sollwert} Die Anpassung des Motors an den eingestellten Sollwert erfolgt, wie erwartet, schneller, wenn die P-Regelung aktiviert ist. Wenn sie deaktiviert ist, dann passiert die  Anpassung etwas verzögert, aber trotzdem recht schnell.
Wenn man nun zusätzlich noch I- und D-Anteil hinzuschaltet, wird die Regelung weiter optimiert.

\subsubsection{Externer Sollwert}

Nun steuert ein Funktionsgenerator den Sollwert über ein eingestelltes Rechteck-Signal. Dabei werden nun die Regelungen getrennt für die P-Regelung, die PI-Regelung und für die PD-Regelung untersucht und schließlich mit der PID-Regelung optimiert.
\paragraph{P-Regelung}
Hier ist der I- und der P-Anteil ausgeschaltet.
\begin{figure}[htbp]
    \begin{subfigure}[b]{0.48\textwidth}
    \centering
    \includegraphics[width=\textwidth]{p_0_skt.jpg}
    \caption{$K_p=0$ Skt}
    \label{fig:p=0}
    \end{subfigure}
    \hfill
    \begin{subfigure}[b]{0.48\textwidth}
    \centering
    \includegraphics[width=\textwidth]{p_10_skt.jpg}
    \caption{$K_p=10$ Skt}
    \label{fig:p=10}
    \end{subfigure}

    \vspace{0.5cm}

    \begin{subfigure}[b]{0.48\textwidth}
    \centering
    \includegraphics[width=\textwidth]{p_50_skt.jpg}
    \caption{$K_p=50$ Skt}
    \label{fig:p=50}
    \end{subfigure}
    \hfill
    \begin{subfigure}[b]{0.48\textwidth}
    \centering
    \includegraphics[width=\textwidth]{p_100_skt.jpg}
    \caption{$K_p=100$ Skt}
    \label{fig:p=100}
    \end{subfigure}

    \vspace{0.5cm}

    \begin{subfigure}[b]{0.48\textwidth}
    \centering
    \includegraphics[width=\textwidth]{p_300_skt.jpg}
    \caption{$K_p=300$ Skt}
    \label{fig:p=300}
    \end{subfigure}
    \hfill
    \begin{subfigure}[b]{0.48\textwidth}
    \centering
    \includegraphics[width=\textwidth]{p_klein_5_skt.jpg}
    \caption{$K_p=5$ Skt, in Blau: Ist-Soll}
    \label{fig:p=5_ist-soll}
    \end{subfigure}
    \caption{P-Regler mit unterschiedlichen $K_p$, in (f) ist zusätzlich ist-soll abgebildet für $K_p=5$ Skt}
    \label{fig:K_p-Regler}
\end{figure}

Man erkennt in \Cref{fig:K_p-Regler} die Anpassung des Motors an das Rechtecksignal (in Blau, Rechtecksignal in Gelb) für unterschiedliche Proportionalbeiwerte $K_p$. In \Cref{fig:p=0} sieht man die Einstellung mit minimalem $K_p$. Dabei ist die Motordrehzahl stets kleiner als das Soll. Bei kleinen \(K_p\) erkennt man, dass das Soll bei der kleineren Amplitude auf Null sinkt und somit der Motor stehen bleibt. Beim Sprung im Rechtecksignal auf die höhere Amplitude wird auch der 
Motor hochgeregelt. Dabei gibt es zunächst einen Peak, der immer höher ist, als das anschließend erreichte Niveau der P-Regelung. Diese Überschwingungen werden größer, mit zunehmendem \(K_p\). Außerdem werden bei steigendem \(K_p\) auch mehr Amplituden der Überschwingung sichtbar. Folglich wird ein stabiler Wert erst später erreicht. Die stationäre Genauigkeit betreffend erkennt man in den Aufnahmen gut, dass diese bei steigendem $K_p$ verbessert wird. Genau erreicht wird das Niveau jedoch nicht, wie man auch bei hohen $K_p$ erkennen kann, aber es nähert sich schon gut an.
Keinen Unterschied sieht man bei der Regelgeschwindigkeit, also der Zeitdauer nach einer Sollwert-Änderung bis zu einer Regel-Anpassung durch den P-Regler.
In \Cref{fig:p=5_ist-soll} erkennt man zusätzlich noch den Ist-Soll-Verlauf. Hierbei wurde die Null auf die Mitte der Anzeige verschoben. Man sieht, dass die Abweichung direkt nach der Sollwert-Anpassung maximal wird, weil der P-Regler natürlich nicht instantan regelt. Ansonsten lässt sich dem Oszilloskop-Bild entnehmen, dass der Soll-Wert sonst immer größer als der Ist-Wert ist.

\paragraph{PI-Regelung} Hier ist nun der P- und der I-Anteil aktiviert, der D-Anteil bleibt weiterhin ausgeschaltet.

\begin{figure}[htbp]
    \begin{subfigure}[b]{0.48\textwidth}
    \centering
    \includegraphics[width=\textwidth]{P_I_0.jpg}
    \caption{$K_p=0$ Skt, $K_I=0$ Skt}
    \label{fig:p=0_i=0}
    \end{subfigure}
    \hfill
    \begin{subfigure}[b]{0.48\textwidth}
    \centering
    \includegraphics[width=\textwidth]{P_0_I_5.jpg}
    \caption{$K_p=0$ Skt, $K_I=5$ Skt}
    \label{fig:p=0_i=5}
    \end{subfigure}
    \vspace{0.5cm}
    \begin{subfigure}[b]{0.48\textwidth}
    \centering
    \includegraphics[width=\textwidth]{P_0_I_50.jpg}
    \caption{$K_p=0$ Skt, $K_I=50$ Skt}
    \label{fig:p=0_i=50}
    \end{subfigure}
    \hfill
    \begin{subfigure}[b]{0.48\textwidth}
    \centering
    \includegraphics[width=\textwidth]{P_50_I_5.jpg}
    \caption{$K_p=50$ Skt, $K_I=5$ Skt}
    \label{fig:p=50_i=5}
    \end{subfigure}
    \vspace{0.5cm}
    \begin{subfigure}[b]{0.48\textwidth}
    \centering
    \includegraphics[width=\textwidth]{P_100_I_50}
    \caption{$K_p=100$ Skt, $K_I=50$ Skt}
    \label{fig:p=100_i=50}
    \end{subfigure}
    \hfill
    \begin{subfigure}[b]{0.48\textwidth}
    \centering
    \includegraphics[width=\textwidth]{P_100_I_10_ist-soll}
    \caption{$K_p=100$ Skt, $K_I=10$ Skt, blau: Ist-Soll}
    \label{fig:p=100_i=50_ist-soll}
    \end{subfigure}
    \caption{PI-Regler mit unterschiedlichen $K_P$ und $K_I$ Kombinationen, in (f): Ist-Soll}
    \label{fig:PI-Regler}
\end{figure}